\section{Circuit Complexity --- Lupanov Gates}
%%% generated by the blueprint pipeline from TCSlib.Complexity.CircuitComplexity.LupanovGates
\subsection*{Circuit Complexity --- Lupanov Gates: Introduction}

Declarations of the Lean module \texttt{TCSlib.Complexity.CircuitComplexity.LupanovGates}.

\subsection*{Circuit Complexity --- Lupanov Gates: Results}

\begin{definition}[Zero-extended Boolean input]
\lean{BoolCircuit.Lupanov.xget}
\label{BoolCircuit.Lupanov.xget}
\leanok
Let $n$ be a natural number, let $x$ be a Boolean string indexed by $0,\ldots,n-1$, and
let $i$ be a natural number. The value at index $i$ is $x_i$ when $i<n$, and is false
when $i\ge n$.
\end{definition}

\begin{definition}[Heap path index]
\lean{BoolCircuit.Lupanov.heapPath}
\label{BoolCircuit.Lupanov.heapPath}
\leanok
\uses{BoolCircuit.Lupanov.xget}
Let $n$ be a natural number, let $x$ be a Boolean string indexed by $0,\ldots,n-1$, and
let $o$ be a natural number. Define $h(0)=1$ and, for every natural number $\ell$,
define $h(\ell+1)=2h(\ell)+b_{o+\ell}$, where $b_j=x_j$ if $j<n$ and $b_j=0$ otherwise.
\end{definition}

\begin{theorem}[Bounds for the heap path value]
\lean{BoolCircuit.Lupanov.heapPath_bounds}
\label{BoolCircuit.Lupanov.heapPath_bounds}
\leanok
\uses{BoolCircuit.Lupanov.heapPath, BoolCircuit.Lupanov.xget}
\difficulty{4}
Let $n$ be a natural number, let $x$ be a Boolean vector indexed by $0,\ldots,n-1$, and
let $o$ be a natural number. For each natural number $l$, form the heap path value by
starting at $1$ and, at each step $j<l$, doubling the current value and adding the
Boolean coordinate of $x$ at index $o+j$ (taking that coordinate to be false when
$o+j\ge n$). Then this value lies in the interval $[2^l,2^{l+1})$.
\end{theorem}

\begin{theorem}[Heap path logarithm]
\lean{BoolCircuit.Lupanov.log_heapPath}
\label{BoolCircuit.Lupanov.log_heapPath}
\leanok
\uses{BoolCircuit.Lupanov.heapPath, BoolCircuit.Lupanov.heapPath_bounds}
\difficulty{1}
Let $n,l,o\in\mathbb{N}$ and let $x:\{0,1\}^n\to\{0,1\}$. The base-$2$ integer logarithm
of the heap path value at depth $l$ with offset $o$ is $l$.
\end{theorem}

\begin{theorem}[Heap-path bits]
\lean{BoolCircuit.Lupanov.testBit_heapPath}
\label{BoolCircuit.Lupanov.testBit_heapPath}
\leanok
\uses{BoolCircuit.Lupanov.heapPath, BoolCircuit.Lupanov.xget}
\difficulty{5}
Let $n,o$ be natural numbers and let $x$ assign a Boolean value to each index below $n$,
extended by false at indices at least $n$. Define $H(0)=1$ and $H(l+1)=2H(l)+x(o+l)$,
where Boolean values are identified with $0$ and $1$. For every $l,i$ with $i<l$, the
binary digit of $H(l)$ in position $l-1-i$ equals $x(o+i)$.
\end{theorem}

\begin{theorem}[Binary digits after adding a power of two]
\lean{BoolCircuit.Lupanov.testBit_two_pow_add_of_lt}
\label{BoolCircuit.Lupanov.testBit_two_pow_add_of_lt}
\leanok
\difficulty{4}
Let $L,t,r$ be natural numbers with $t<2^L$. The binary digit in position $r$ of $2^L+t$
is true exactly when either the binary digit of $t$ in position $r$ is true or $r=L$.
\end{theorem}

\begin{definition}[Integer encoded by a bit list]
\lean{BoolCircuit.Lupanov.ofBits}
\label{BoolCircuit.Lupanov.ofBits}
\leanok
For a finite list of Boolean values, define its encoded natural number recursively: the
empty list encodes $0$, and a list with first bit $b$ and remaining list $L$ encodes
twice the number encoded by $L$, plus $1$ if $b$ is true and $0$ if $b$ is false.
\end{definition}

\begin{theorem}[Bits of a Boolean list]
\lean{BoolCircuit.Lupanov.testBit_ofBits}
\label{BoolCircuit.Lupanov.testBit_ofBits}
\leanok
\uses{BoolCircuit.Lupanov.ofBits}
\difficulty{4}
Encode a finite Boolean list recursively by assigning the empty list the integer $0$ and
assigning a list with first entry $b$ and remaining list $L$ the integer $2E(L)+b$,
where $E$ denotes this encoding and Boolean values are identified with $0$ and $1$. For
every list $L$ and index $i$, the binary digit in position $i$ of its encoding equals
the entry at index $i$, taking the entry to be false when the index is outside the list.
\end{theorem}

\begin{theorem}[Bound on the value of a bit list]
\lean{BoolCircuit.Lupanov.ofBits_lt}
\label{BoolCircuit.Lupanov.ofBits_lt}
\leanok
\uses{BoolCircuit.Lupanov.ofBits}
\difficulty{4}
For every finite list $l$ of Boolean bits, its binary value is strictly less than
$2^{|l|}$, where $|l|$ is the length of the list.
\end{theorem}

\begin{definition}[Literal vertex index]
\lean{BoolCircuit.Lupanov.litV}
\label{BoolCircuit.Lupanov.litV}
\leanok
For natural numbers $n$ and $i$ and a Boolean value $b$, assign the vertex index $i$
when $b$ is true and the vertex index $n+i$ when $b$ is false.
\end{definition}

\begin{definition}[Heap gate]
\lean{BoolCircuit.Lupanov.heapGate}
\label{BoolCircuit.Lupanov.heapGate}
\leanok
\uses{BoolCircuit.DAGGate, BoolCircuit.Lupanov.litV}
For natural numbers $n$, $\mathrm{base}$, $o$, and $q$, define a gate of type AND as
follows. If $q\le 1$, it has no input vertices; otherwise its input vertices are
$\mathrm{base}+\max(\lfloor q/2\rfloor-1,0)$ and the literal vertex with index
$o+\max(\lfloor\log_2 q\rfloor-1,0)$, using $n$ to determine the literal vertex, with
the Boolean label true exactly when $q$ is odd.
\end{definition}

\begin{definition}[Function gate]
\lean{BoolCircuit.Lupanov.funcGate}
\label{BoolCircuit.Lupanov.funcGate}
\leanok
\uses{BoolCircuit.DAGGate}
For natural numbers $f_{\mathrm{base}}$, $z_{\mathrm{base}}$, $k$, and $t$, define a
gate of type OR. If $t=0$, it has no input vertices; otherwise its inputs are
$f_{\mathrm{base}}+\max(t-2^{\lfloor\log_2 t\rfloor},0)$ and
$z_{\mathrm{base}}+\max(2^k+\lfloor\log_2 t\rfloor-1,0)$.
\end{definition}

\begin{definition}[Combination gate]
\lean{BoolCircuit.Lupanov.combGate}
\label{BoolCircuit.Lupanov.combGate}
\leanok
\uses{BoolCircuit.DAGGate}
Let $y_{\mathrm{base}}$, $f_{\mathrm{base}}$, $m$, and $p$ be natural numbers, and let
$c$ be a function from natural numbers to natural numbers. Define a gate of type AND
whose two input vertices are $y_{\mathrm{base}}+\max(2^m+p-1,0)$ and
$f_{\mathrm{base}}+c(p)$.
\end{definition}

\begin{definition}[First gate-count component]
\lean{BoolCircuit.Lupanov.oZ}
\label{BoolCircuit.Lupanov.oZ}
\leanok
For natural numbers $n$ and $k$, define $z(n,k)=n+(2^{n-k+1}-1)$, where subtraction in
the exponent is truncated natural-number subtraction.
\end{definition}

\begin{definition}[Second gate-count component]
\lean{BoolCircuit.Lupanov.oF}
\label{BoolCircuit.Lupanov.oF}
\leanok
\uses{BoolCircuit.Lupanov.oZ}
For natural numbers $n$ and $k$, define $f(n,k)=z(n,k)+(2^{k+1}-1)$, where
$z(n,k)=n+(2^{n-k+1}-1)$ and subtraction is truncated natural-number subtraction.
\end{definition}

\begin{definition}[Third gate-count component]
\lean{BoolCircuit.Lupanov.oC}
\label{BoolCircuit.Lupanov.oC}
\leanok
\uses{BoolCircuit.Lupanov.oF}
For natural numbers $n$ and $k$, define $c(n,k)=f(n,k)+2^{2^k}$, where
$f(n,k)=z(n,k)+(2^{k+1}-1)$ and $z(n,k)=n+(2^{n-k+1}-1)$; subtraction in the exponent
defining $z$ is truncated natural-number subtraction.
\end{definition}

\begin{definition}[Total gate-count bound]
\lean{BoolCircuit.Lupanov.total}
\label{BoolCircuit.Lupanov.total}
\leanok
\uses{BoolCircuit.Lupanov.oC}
For natural numbers $n$ and $k$, define $T(n,k)=c(n,k)+2^{n-k}$, where subtraction is
truncated natural-number subtraction and $c(n,k)=f(n,k)+2^{2^k}$, with
$f(n,k)=z(n,k)+(2^{k+1}-1)$ and $z(n,k)=n+(2^{n-k+1}-1)$.
\end{definition}

\begin{definition}[Input and block assembly]
\lean{BoolCircuit.Lupanov.assemble}
\label{BoolCircuit.Lupanov.assemble}
\leanok
Let $n,k,p,r$ be natural numbers. Define a Boolean string $x$ of length $n$ by setting,
for each $i<n$, $x_i$ to the bit in position $n-k-1-i$ of $2^{n-k}+p$ when $i<n-k$, and
to the bit in position $k-1-(i-(n-k))$ of $2^k+r$ otherwise; all subtractions here are
truncated natural-number subtraction.
\end{definition}

\begin{definition}[Truth-table block code]
\lean{BoolCircuit.Lupanov.code}
\label{BoolCircuit.Lupanov.code}
\leanok
\uses{BoolCircuit.Lupanov.assemble, BoolCircuit.Lupanov.ofBits}
Let $n$ be a natural number, let $f$ be a Boolean function on strings of length $n$, and
let $k,p$ be natural numbers. Define the code to be the natural number obtained by
interpreting as bits the list of values $f(x_r)$, indexed by $0\le r<2^k$, where $x_r$
is the string assembled from $n,k,p,r$.
\end{definition}

\begin{definition}[Gate at a specified index]
\lean{BoolCircuit.Lupanov.gateAt}
\label{BoolCircuit.Lupanov.gateAt}
\leanok
\uses{BoolCircuit.DAGGate, BoolCircuit.Lupanov.code, BoolCircuit.Lupanov.combGate, BoolCircuit.Lupanov.funcGate, BoolCircuit.Lupanov.heapGate, BoolCircuit.Lupanov.oC, BoolCircuit.Lupanov.oF, BoolCircuit.Lupanov.oZ}
Let $n,k,j$ be natural numbers and let $f$ be a Boolean function on strings of length
$n$. The gate at index $j$ is a NOT gate reading input $j$ if $j<n$; otherwise it is
selected from the heap-gate range when $j<z(n,k)$, from the function-gate range when
$j<f(n,k)$, from the function-code gate range when $j<c(n,k)$, and from the
combination-gate range thereafter, with the corresponding offsets and parameters
determined by these ranges.
\end{definition}

\begin{definition}[List of circuit gates]
\lean{BoolCircuit.Lupanov.gates}
\label{BoolCircuit.Lupanov.gates}
\leanok
\uses{BoolCircuit.DAGGate, BoolCircuit.Lupanov.gateAt, BoolCircuit.Lupanov.total}
Let $n,k$ be natural numbers and let $f$ be a Boolean function on strings of length $n$.
Define the gate list to consist, in order, of the gates at indices $0$ through
$T(n,k)-1$, where $T(n,k)=c(n,k)+2^{n-k}$, $c(n,k)=f_0(n,k)+2^{2^k}$,
$f_0(n,k)=z(n,k)+(2^{k+1}-1)$, and $z(n,k)=n+(2^{n-k+1}-1)$; all natural-number
subtractions are truncated.
\end{definition}

\begin{definition}[Comb root indices]
\lean{BoolCircuit.Lupanov.combRoots}
\label{BoolCircuit.Lupanov.combRoots}
\leanok
\uses{BoolCircuit.Lupanov.oC}
For natural numbers $n$ and $k$, the comb roots are the list of indices $n + o_C(n,k) +
p$ as $p$ ranges from $0$ to $2^{n-k}-1$.
\end{definition}

\begin{definition}[Lupanov gate output]
\lean{BoolCircuit.Lupanov.lupGates}
\label{BoolCircuit.Lupanov.lupGates}
\leanok
\uses{BoolCircuit.DAGGate, BoolCircuit.Lupanov.combRoots, BoolCircuit.Lupanov.gates, BoolCircuit.emitTree}
Given a Boolean function $f : \{0,1\}^n \to \{0,1\}$ and a natural number $k$, the
Lupanov gate output is the pair obtained by emitting an OR tree whose roots are the comb
roots for $n$ and $k$, using the gate list associated with $f$ and $k$.
\end{definition}

\begin{theorem}[Length of the gate list]
\lean{BoolCircuit.Lupanov.length_gates}
\label{BoolCircuit.Lupanov.length_gates}
\leanok
\uses{BoolCircuit.Lupanov.gates, BoolCircuit.Lupanov.total}
\difficulty{1}
Let $n,k\in\mathbb{N}$ and let $f:\{0,1\}^n\to\{0,1\}$. The length of the gate list
constructed from $f$ and $k$ is $\mathrm{total}(n,k)$.
\end{theorem}

\begin{theorem}[Gate list entry]
\lean{BoolCircuit.Lupanov.getElem_gates}
\label{BoolCircuit.Lupanov.getElem_gates}
\leanok
\uses{BoolCircuit.Lupanov.gateAt, BoolCircuit.Lupanov.gates}
\difficulty{1}
Let $n,k,j\in\mathbb{N}$ and let $f:\{0,1\}^n\to\{0,1\}$. Whenever $j$ is less than the
length of the constructed gate list for $f$ and $k$, its entry at index $j$ is the gate
assigned at index $j$ by the gate construction.
\end{theorem}

\begin{theorem}[Successor power of two]
\lean{BoolCircuit.Lupanov.two_pow_succ_eq}
\label{BoolCircuit.Lupanov.two_pow_succ_eq}
\leanok
\difficulty{3}
For every natural number $a$, $2^{a+1}=2\cdot2^a$.
\end{theorem}

\begin{theorem}[Literal vertex bound]
\lean{BoolCircuit.Lupanov.litV_le}
\label{BoolCircuit.Lupanov.litV_le}
\leanok
\uses{BoolCircuit.Lupanov.litV}
\difficulty{4}
For every $n,i\in\mathbb{N}$ and every Boolean value $b$, the literal vertex associated
to $n,i,b$ is at most $n+i$.
\end{theorem}

\begin{theorem}[A logarithmic bound from a power bound]
\lean{BoolCircuit.Lupanov.log_le_of_lt_two_pow_succ}
\label{BoolCircuit.Lupanov.log_le_of_lt_two_pow_succ}
\leanok
\difficulty{4}
For natural numbers $q$ and $a$, if $q<2^{a+1}$, then the base-$2$ logarithm of $q$ is
at most $a$.
\end{theorem}

\begin{theorem}[An empty-input gate has fan-in at most two]
\lean{BoolCircuit.Lupanov.nil_faninTwo}
\label{BoolCircuit.Lupanov.nil_faninTwo}
\leanok
\uses{BoolCircuit.DAGGate, BoolCircuit.DAGGate.FaninTwo, BoolCircuit.GateKind}
\difficulty{3}
Let $k$ be an AND, OR, or NOT gate kind other than NOT. The gate of kind $k$ with no
input vertices is well formed and has at most two inputs.
\end{theorem}

\begin{theorem}[Two-input gate well-formedness]
\lean{BoolCircuit.Lupanov.pair_faninTwo}
\label{BoolCircuit.Lupanov.pair_faninTwo}
\leanok
\uses{BoolCircuit.DAGGate, BoolCircuit.DAGGate.FaninTwo, BoolCircuit.GateKind}
\difficulty{1}
Let $k$ be an AND, OR, or NOT gate kind with $k\ne\mathrm{NOT}$, and let $a,b$ be
distinct natural numbers. The gate of kind $k$ with inputs $a$ and $b$ is well formed
and has at most two inputs.
\end{theorem}

\begin{theorem}[Gate inputs precede their gate]
\lean{BoolCircuit.Lupanov.gateAt_args_lt}
\label{BoolCircuit.Lupanov.gateAt_args_lt}
\leanok
\uses{BoolCircuit.Lupanov.assemble, BoolCircuit.Lupanov.code, BoolCircuit.Lupanov.combGate, BoolCircuit.Lupanov.funcGate, BoolCircuit.Lupanov.gateAt, BoolCircuit.Lupanov.heapGate, BoolCircuit.Lupanov.litV_le, BoolCircuit.Lupanov.log_le_of_lt_two_pow_succ, BoolCircuit.Lupanov.oC, BoolCircuit.Lupanov.oF, BoolCircuit.Lupanov.oZ, BoolCircuit.Lupanov.ofBits_lt, BoolCircuit.Lupanov.total, BoolCircuit.Lupanov.two_pow_succ_eq}
\difficulty{6}
Let $n,k\in\mathbb{N}$ with $k\le n$, let $f:\{0,1\}^n\to\{0,1\}$, and let $j$ be a gate
index below the total number of gates in the constructed circuit. Every vertex read by
the gate at index $j$ has index strictly less than $n+j$.
\end{theorem}

\begin{theorem}[Circuit gates have fan-in at most two]
\lean{BoolCircuit.Lupanov.gateAt_faninTwo}
\label{BoolCircuit.Lupanov.gateAt_faninTwo}
\leanok
\uses{BoolCircuit.DAGGate.FaninTwo, BoolCircuit.Lupanov.funcGate, BoolCircuit.Lupanov.gateAt, BoolCircuit.Lupanov.heapGate, BoolCircuit.Lupanov.litV_le, BoolCircuit.Lupanov.log_le_of_lt_two_pow_succ, BoolCircuit.Lupanov.nil_faninTwo, BoolCircuit.Lupanov.oC, BoolCircuit.Lupanov.oF, BoolCircuit.Lupanov.oZ, BoolCircuit.Lupanov.pair_faninTwo, BoolCircuit.Lupanov.total, BoolCircuit.Lupanov.two_pow_succ_eq}
\difficulty{5}
Let $n,k\in\mathbb{N}$ with $k\le n$, let $f:\{0,1\}^n\to\{0,1\}$, and let $j$ be a gate
index below the total number of gates in the constructed circuit. The gate at index $j$
is well formed and reads at most two vertices.
\end{theorem}

\begin{theorem}[Acyclicity of the gate list]
\lean{BoolCircuit.Lupanov.gates_acyclic}
\label{BoolCircuit.Lupanov.gates_acyclic}
\leanok
\uses{BoolCircuit.GatesAcyclic, BoolCircuit.Lupanov.gateAt_args_lt, BoolCircuit.Lupanov.gates, BoolCircuit.Lupanov.getElem_gates}
\difficulty{3}
Let $n,k\in\mathbb{N}$ with $k\le n$, and let $f:\{0,1\}^n\to\{0,1\}$. In the gate list
constructed from $f$ and $k$, every gate at index $i$ has all its input vertices indexed
below $n+i$.
\end{theorem}

\begin{theorem}[Lupanov offset identities]
\lean{BoolCircuit.Lupanov.offsets}
\label{BoolCircuit.Lupanov.offsets}
\leanok
\uses{BoolCircuit.Lupanov.oC, BoolCircuit.Lupanov.oF, BoolCircuit.Lupanov.oZ, BoolCircuit.Lupanov.total, BoolCircuit.Lupanov.two_pow_succ_eq}
\difficulty{3}
For natural numbers $n$ and $k$, let $o_Z,o_F,o_C$, and $T$ denote the offsets defined
by $o_Z=n+(2^{n-k+1}-1)$, $o_F=o_Z+(2^{k+1}-1)$, $o_C=o_F+2^{2^k}$, and $T=o_C+2^{n-k}$.
Then $o_Z+1=n+2\cdot2^{n-k}$, $o_F+1=o_Z+2\cdot2^k$, and the stated formulas for $o_C$
and $T$ hold; moreover, $1\le 2^{n-k}$, $1\le 2^k$, and $1\le 2^{2^k}$.
\end{theorem}

\begin{theorem}[First heap block gate]
\lean{BoolCircuit.Lupanov.gateAt_y}
\label{BoolCircuit.Lupanov.gateAt_y}
\leanok
\uses{BoolCircuit.Lupanov.gateAt, BoolCircuit.Lupanov.heapGate, BoolCircuit.Lupanov.oZ, BoolCircuit.Lupanov.offsets, BoolCircuit.Lupanov.two_pow_succ_eq}
\difficulty{4}
Let $n,k,q\in\mathbb{N}$ and let $f:\{0,1\}^n\to\{0,1\}$. If $1\le q<2^{n-k+1}$, then
the gate at index $n+(q-1)$ in the construction for $f$ and parameter $k$ is the heap
gate with parameters $n$, base $2n$, offset $0$, and size $q$.
\end{theorem}

\begin{theorem}[Second heap block gate]
\lean{BoolCircuit.Lupanov.gateAt_z}
\label{BoolCircuit.Lupanov.gateAt_z}
\leanok
\uses{BoolCircuit.Lupanov.gateAt, BoolCircuit.Lupanov.heapGate, BoolCircuit.Lupanov.oF, BoolCircuit.Lupanov.oZ, BoolCircuit.Lupanov.offsets, BoolCircuit.Lupanov.two_pow_succ_eq}
\difficulty{4}
Let $n,k,q\in\mathbb{N}$ and let $f:\{0,1\}^n\to\{0,1\}$. If $1\le q<2^{k+1}$, then the
gate at index $o_Z(n,k)+(q-1)$ in the construction for $f$ and parameter $k$ is the heap
gate with parameters $n$, base $n+o_Z(n,k)$, offset $n-k$, and size $q$.
\end{theorem}

\begin{theorem}[Function gate in the circuit]
\lean{BoolCircuit.Lupanov.gateAt_f}
\label{BoolCircuit.Lupanov.gateAt_f}
\leanok
\uses{BoolCircuit.Lupanov.funcGate, BoolCircuit.Lupanov.gateAt, BoolCircuit.Lupanov.oC, BoolCircuit.Lupanov.oF, BoolCircuit.Lupanov.oZ, BoolCircuit.Lupanov.offsets}
\difficulty{4}
Let $f:\{0,1\}^n\to\{0,1\}$ and $k,t\in\mathbb{N}$, with $t<2^{2^k}$. The gate at index
$o_F(n,k)+t$ is the function gate with bases $n+o_F(n,k)$ and $n+o_Z(n,k)$, parameter
$k$, and index parameter $t$.
\end{theorem}

\begin{theorem}[Combination gate in the circuit]
\lean{BoolCircuit.Lupanov.gateAt_c}
\label{BoolCircuit.Lupanov.gateAt_c}
\leanok
\uses{BoolCircuit.Lupanov.code, BoolCircuit.Lupanov.combGate, BoolCircuit.Lupanov.gateAt, BoolCircuit.Lupanov.oC, BoolCircuit.Lupanov.oF, BoolCircuit.Lupanov.oZ, BoolCircuit.Lupanov.offsets}
\difficulty{4}
For any Boolean function $f:\{0,1\}^n\to\{0,1\}$ and natural numbers $k,p$, the gate at
index $o_C(n,k)+p$ is the combination gate with bases $n+n$ and $n+o_F(n,k)$, parameter
$n-k$, code given by $f$ and $k$, and index parameter $p$.
\end{theorem}

\begin{theorem}[Value at a circuit gate]
\lean{BoolCircuit.Lupanov.vertexValue_gate}
\label{BoolCircuit.Lupanov.vertexValue_gate}
\leanok
\uses{BoolCircuit.DAGGate, BoolCircuit.DAGGate.eval, BoolCircuit.DAGGate.eval_congr, BoolCircuit.runWith, BoolCircuit.runWith_getD_gate, BoolCircuit.runWith_getD_take, BoolCircuit.vertexValue}
\difficulty{4}
Let $G$ be a finite list of Boolean gates and let $x$ be a Boolean input indexed by
$0,\ldots,n-1$. If $j$ is less than the length of $G$ and every argument index of the
$j$th gate is less than $n+j$, then the value at index $n+j$ in the list obtained by
starting with $x$ and successively appending each gate’s value equals the evaluation of
the $j$th gate on that resulting list.
\end{theorem}

\begin{theorem}[Value of a circuit gate]
\lean{BoolCircuit.Lupanov.value_gateAt}
\label{BoolCircuit.Lupanov.value_gateAt}
\leanok
\uses{BoolCircuit.DAGGate.eval, BoolCircuit.Lupanov.gateAt, BoolCircuit.Lupanov.gateAt_args_lt, BoolCircuit.Lupanov.gates, BoolCircuit.Lupanov.getElem_gates, BoolCircuit.Lupanov.total, BoolCircuit.Lupanov.vertexValue_gate, BoolCircuit.runWith, BoolCircuit.vertexValue}
\difficulty{3}
Let $n,k$ be natural numbers with $k\le n$, let $f:\{0,1\}^n\to\{0,1\}$ be a Boolean
function, and let $x\in\{0,1\}^n$. For every natural number $j<\mathrm{total}(n,k)$, the
value at vertex $n+j$ after evaluating the circuit formed by the gates is the evaluation
of the gate at position $j$ on the resulting list of vertex values.
\end{theorem}

\begin{theorem}[Evaluation of a two-input conjunction gate]
\lean{BoolCircuit.Lupanov.eval_and_pair}
\label{BoolCircuit.Lupanov.eval_and_pair}
\leanok
\uses{BoolCircuit.DAGGate, BoolCircuit.DAGGate.eval}
\difficulty{1}
For any natural numbers $a,b$ and any Boolean value list $v$, the value of a conjunction
gate reading positions $a$ and $b$ is $v_a\land v_b$, where an out-of-range list
position is interpreted as false.
\end{theorem}

\begin{theorem}[Evaluation of a two-input disjunction gate]
\lean{BoolCircuit.Lupanov.eval_or_pair}
\label{BoolCircuit.Lupanov.eval_or_pair}
\leanok
\uses{BoolCircuit.DAGGate, BoolCircuit.DAGGate.eval}
\difficulty{1}
For any natural numbers $a,b$ and any Boolean value list $v$, the value of a disjunction
gate reading positions $a$ and $b$ is $v_a\lor v_b$, where an out-of-range list position
is interpreted as false.
\end{theorem}

\begin{theorem}[Values of literal vertices]
\lean{BoolCircuit.Lupanov.value_litV}
\label{BoolCircuit.Lupanov.value_litV}
\leanok
\uses{BoolCircuit.DAGGate.eval, BoolCircuit.Lupanov.gateAt, BoolCircuit.Lupanov.gates, BoolCircuit.Lupanov.litV, BoolCircuit.Lupanov.offsets, BoolCircuit.Lupanov.value_gateAt, BoolCircuit.Lupanov.xget, BoolCircuit.vertexValue, BoolCircuit.vertexValue_input}
\difficulty{5}
Let $n,k$ be natural numbers with $k\le n$, let $f:\{0,1\}^n\to\{0,1\}$ and
$x\in\{0,1\}^n$, and take $i<n$ and $b\in\{0,1\}$. The value at vertex $i$ if $b$ is
true, and at vertex $n+i$ if $b$ is false, equals true exactly when the input bit $x_i$
equals $b$.
\end{theorem}

\begin{theorem}[Values of heap gates]
\lean{BoolCircuit.Lupanov.value_heap}
\label{BoolCircuit.Lupanov.value_heap}
\leanok
\uses{BoolCircuit.DAGGate.eval, BoolCircuit.Lupanov.eval_and_pair, BoolCircuit.Lupanov.gateAt, BoolCircuit.Lupanov.gates, BoolCircuit.Lupanov.heapGate, BoolCircuit.Lupanov.heapPath, BoolCircuit.Lupanov.litV, BoolCircuit.Lupanov.log_le_of_lt_two_pow_succ, BoolCircuit.Lupanov.total, BoolCircuit.Lupanov.value_gateAt, BoolCircuit.Lupanov.value_litV, BoolCircuit.Lupanov.xget, BoolCircuit.vertexValue}
\difficulty{6}
Let $n,k$ be natural numbers with $k\le n$, let $f:\{0,1\}^n\to\{0,1\}$ and
$x\in\{0,1\}^n$, and let $bj,o,L$ be natural numbers with $o+L\le n$. Suppose that for
every $q$ with $1\le q<2^{L+1}$, the gate at position $bj+(q-1)$ is the heap gate with
parameters $n$, base $n+bj$, offset $o$, and index $q$. Also suppose
$bj+(2^{L+1}-1)\le\mathrm{total}(n,k)$. Then for every $q$ with $1\le q<2^{L+1}$, the
value at vertex $n+bj+(q-1)$ is true exactly when the heap path of $x$ from offset $o$
at depth $\lfloor\log_2 q\rfloor$ equals $q$.
\end{theorem}

\begin{theorem}[Logarithm of a power plus a remainder]
\lean{BoolCircuit.Lupanov.log_two_pow_add}
\label{BoolCircuit.Lupanov.log_two_pow_add}
\leanok
\uses{BoolCircuit.Lupanov.two_pow_succ_eq}
\difficulty{3}
For natural numbers $a$ and $p$, if $p<2^a$, then the base-$2$ logarithm of $2^a+p$
equals $a$.
\end{theorem}

\begin{theorem}[Values of the first-level leaves]
\lean{BoolCircuit.Lupanov.value_yleaf}
\label{BoolCircuit.Lupanov.value_yleaf}
\leanok
\uses{BoolCircuit.Lupanov.gateAt_y, BoolCircuit.Lupanov.gates, BoolCircuit.Lupanov.heapPath, BoolCircuit.Lupanov.log_two_pow_add, BoolCircuit.Lupanov.offsets, BoolCircuit.Lupanov.two_pow_succ_eq, BoolCircuit.Lupanov.value_heap, BoolCircuit.vertexValue}
\difficulty{3}
Let $n,k$ be natural numbers with $k\le n$, let $f:\{0,1\}^n\to\{0,1\}$ and
$x\in\{0,1\}^n$, and take $p<2^{n-k}$. The value at vertex $2n+(2^{n-k}+p-1)$ is true
exactly when the heap path of $x$ from offset $0$ at depth $n-k$ equals $2^{n-k}+p$.
\end{theorem}

\begin{theorem}[Values of the second-level leaves]
\lean{BoolCircuit.Lupanov.value_zleaf}
\label{BoolCircuit.Lupanov.value_zleaf}
\leanok
\uses{BoolCircuit.Lupanov.gateAt_z, BoolCircuit.Lupanov.gates, BoolCircuit.Lupanov.heapPath, BoolCircuit.Lupanov.log_two_pow_add, BoolCircuit.Lupanov.oZ, BoolCircuit.Lupanov.offsets, BoolCircuit.Lupanov.two_pow_succ_eq, BoolCircuit.Lupanov.value_heap, BoolCircuit.vertexValue}
\difficulty{3}
Let $n,k$ be natural numbers with $k\le n$, let $f:\{0,1\}^n\to\{0,1\}$ and
$x\in\{0,1\}^n$, and take $r<2^k$. The value at vertex $n+o_Z(n,k)+(2^k+r-1)$ is true
exactly when the heap path of $x$ from offset $n-k$ at depth $k$ equals $2^k+r$.
\end{theorem}

\begin{theorem}[Values of function gates]
\lean{BoolCircuit.Lupanov.value_func}
\label{BoolCircuit.Lupanov.value_func}
\leanok
\uses{BoolCircuit.DAGGate.eval, BoolCircuit.Lupanov.eval_or_pair, BoolCircuit.Lupanov.funcGate, BoolCircuit.Lupanov.gateAt_f, BoolCircuit.Lupanov.gates, BoolCircuit.Lupanov.heapPath, BoolCircuit.Lupanov.heapPath_bounds, BoolCircuit.Lupanov.oF, BoolCircuit.Lupanov.oZ, BoolCircuit.Lupanov.offsets, BoolCircuit.Lupanov.testBit_two_pow_add_of_lt, BoolCircuit.Lupanov.two_pow_succ_eq, BoolCircuit.Lupanov.value_gateAt, BoolCircuit.Lupanov.value_zleaf, BoolCircuit.vertexValue}
\difficulty{6}
Let $n$ be a natural number, let $f$ be a Boolean function on inputs indexed below $n$,
and let $k\le n$ be a natural number. For every Boolean input $x$ on those indices and
every $t<2^{2^k}$, the value at vertex $n+o_F(n,k)+t$ of the gate list constructed from
$f$ and $k$ equals the binary digit at position $H_x(n-k,k)-2^k$ of $t$, where
$H_x(o,0)=1$ and $H_x(o,l+1)=2H_x(o,l)+x(o+l)$, with $x$ taken to be false outside its
input indices.
\end{theorem}

\begin{theorem}[Value of a combining gate]
\lean{BoolCircuit.Lupanov.value_comb}
\label{BoolCircuit.Lupanov.value_comb}
\leanok
\uses{BoolCircuit.Lupanov.assemble, BoolCircuit.Lupanov.code, BoolCircuit.Lupanov.combGate, BoolCircuit.Lupanov.eval_and_pair, BoolCircuit.Lupanov.gateAt_c, BoolCircuit.Lupanov.gates, BoolCircuit.Lupanov.heapPath, BoolCircuit.Lupanov.oC, BoolCircuit.Lupanov.oF, BoolCircuit.Lupanov.ofBits_lt, BoolCircuit.Lupanov.offsets, BoolCircuit.Lupanov.value_func, BoolCircuit.Lupanov.value_gateAt, BoolCircuit.Lupanov.value_yleaf, BoolCircuit.vertexValue}
\difficulty{4}
Let $n$ be a natural number, let $f$ be a Boolean function on the Boolean inputs indexed
below $n$, let $k$ be a natural number with $k\le n$, and let $x$ be such an input. For
every natural number $p<2^{n-k}$, the value at vertex $n+o_C(n,k)+p$ of the gate list
constructed from $f$ and $k$ equals the conjunction of the following two Boolean
conditions: the heap-path integer of $x$ at depth $n-k$ equals $2^{n-k}+p$, and the
binary digit at position $H_x(n-k,k)-2^k$ of the code of $f$ for $k,p$ is true. Here
$H_x(o,l)$ is defined by $H_x(o,0)=1$ and $H_x(o,l+1)=2H_x(o,l)+x(o+l)$, with $x$ taken
to be false outside its input indices.
\end{theorem}
